\section{\label{sec:conclusion}Conclusion}

We studied feasible-subspace mixers for a fixed-cardinality binary portfolio
problem, where the constraint confines the dynamics to a Hamming-weight sector
and removes the probability leakage that penalty encodings
incur~\cite{Hadfield2019,Slate2021}.

The theoretical contribution is a correspondence between some types of mixers and quantum walks. Lemma~\ref{lem:dictionary} states
the conditions under which a feasibility-preserving mixer is the adjacency
operator of a graph on the feasible set and identifies that graph, and
Remark~\ref{rem:failure} records what the mixer generates when the conditions
fail: an on-site potential term when the diagonal is not constant, and
dynamics that is not that of any real-weighted walk when the off-diagonal
phases do not cancel around cycles. The correspondence does not address the
construction of subspace-preserving mixers, for which we refer to
Refs.~\cite{Fuchs2022,Fuchs2024}.

Three consequences follow. The complete $XY$ mixer with a Dicke initial state
is exactly QWOA on the Johnson graph $J(n,k)$, with no global phase and no
rescaling. The Grover mixer of Ref.~\cite{Bartschi2020} and the
complete-graph QWOA mixer of Ref.~\cite{Slate2021} are the same operator up to
a global phase and the rescaling $\theta=M\beta/2$; the equivalence is stated
in Ref.~\cite{Bridi2024}. This transfers the
structure-independence result of Ref.~\cite{Bridi2024} from the Grover mixer
to complete-graph QWOA. And the graph on which the ring $XY$ mixer walks has
vertex degree $2b(x)$ and is therefore not regular;
Proposition~\ref{prop:xy-regularity} shows that this is the generic case, the
regular mixers being those of the complete graph together with the star and
its complement at $n=2k$.

These place the two most widely used constrained mixers at opposite ends of
a single family of regular, vertex-transitive graphs on $\F_{n,k}$.
Reference~\cite{Slate2021} argues that the complete graph is likely the
best choice of feasible mixer, by analogy with its optimality for spatial
search~\cite{Zalka1999,Childs2004,Roland2003}. We find the opposite
ordering, and not as a trend in density.

{\color{red}\textbf{[REVIEW -- scope]} Qualify as an empirical result under the present benchmark and optimization protocol, not a universal mixer ordering.}

{\color{red}\textbf{[LITERATURE UPDATE -- revised positioning]} In light of Matwiejew and Wang~\cite{Matwiejew2024}, He et al.~\cite{He2023}, Kordonowy and Leipold~\cite{Kordonowy2026}, and Bridi et al.~\cite{Bridi2026Expressivity}, the strongest defensible message is not that sparsity or Johnson connectivity is universally superior. A stronger literature-aligned conclusion is that feasible-state graph geometry, spectrum, initial-state alignment, expressivity/trainability, and implementation error jointly determine mixer quality; the present Johnson-scheme sweep isolates one controlled family in which complete feasible mixing is a consistently weak endpoint under the tested protocol.}

After a single layer, most unions of distance classes are statistically
indistinguishable from $J(n,k)$, whose degree is smaller by a factor of
three to fourteen depending on the size. A few unions fall between that
cluster and the complete graph, which sits furthest apart at every size and
on all three figures of merit. With depth the separation widens, $J(n,k)$
reaching $P_*=0.040$ against $0.008$ at $(16,4)$ and $p=6$. The complete
graph is thus an outlier within the family rather than one end of a
gradient, and it is the one member that discards all information about the
combinatorial distance between configurations.

Circuit cost sharpens this. Only $J(n,k)$ is two-local; every union
containing a class $j\ge2$ needs $2j$-body generators; and the complete
graph, though exact, costs roughly eight times the \textsc{cnot}s of one
Trotter step of the $XY$ mixer on $K_n$ at every size, and thirty to sixty
times the depth.

{\color{red}\textbf{[REVIEW -- matched-error requirement]} These are raw compilation ratios versus one approximate Trotter step, not equal-accuracy resource ratios.} Constraint preservation and connectivity are separate
design choices, and the second is not settled by maximizing it.

Two caveats on scope. Related Dicke-state portfolio work now also includes the multiclass VQE framework of Scursulim et al.~\cite{Scursulim2026}, which enforces diversification/cardinality constraints in the ansatz and explicitly studies optimizer dependence. This reinforces that the novelty here should be assigned to feasible-state mixer geometry and the Johnson-scheme comparison, not to Dicke-state feasibility encoding itself.

{\color{red}\textbf{[LITERATURE UPDATE -- positioning]} Cite Ref.~\cite{Scursulim2026} when delimiting the portfolio contribution. It is especially useful because its CMA-ES/COBYLA comparison independently supports treating classical optimization as an experimental variable rather than an implementation detail.}

The benchmark is not a reproduction of
Ref.~\cite{Slate2021}, whose three-position encoding gives a different
feasible set, so it tests the general recommendation rather than the
specific result. And the ordering at the endpoints is what
Ref.~\cite{Bridi2024} predicts; what is new is that the intermediate graphs
do not interpolate between them.

Several limitations bound the conclusions. We use a single risk-aversion
parameter, chosen because the fraction of instances surviving the hardness
screen is strongly monotone in it, so the low-risk regime is excluded by
construction. The ordering is not stable in system size: at $(20,5)$ and
$p=4$ the leading mixers are $R_{23}$ and $R_{13}$ rather than $J(n,k)$, by
a factor of about $1.5$ in $P_*$, and six instances cannot settle whether
this is a reversal. Only the position of the complete graph survives every
size tested.

{\color{red}\textbf{[REVIEW -- optimizer caveat]} Add ``under the present optimization protocol,'' because the manuscript itself reports substantial optimizer gaps at $p=1$.} Optimized values at every depth are lower bounds: on common
instances the exhaustive single-layer scan exceeds the optimizer at $p=1$
by a median factor of $1.1$ to $1.8$, so part of any gap between mixers may
reflect the difficulty of locating good parameters rather than their
absence. Finally, every member of the family here is regular and
vertex-transitive, which is what makes the comparison clean and also what
confines it, since Lemma~\ref{lem:dictionary} admits feasible graphs that
are neither. How the effect behaves beyond exact simulation is open, as is
the circuit-level cost of feasible mixers under realistic noise.

\section*{Data and code availability}
% CONFIRMAR: inserir DOI/URL do repositorio

{\color{red}\textbf{[REVIEW -- submission blocker]} Replace with a permanent repository release/DOI, freeze code/data, record dependency versions, and provide a scripted reproduction path.}
The code that generates the instances, runs the simulations and produces
every figure and table of this work, together with the resulting data, is
available at \url{<repository>}.